immediately yields a factor
\[
  (-1)^{(F-1)+(F-2)+\cdots+1}=(-1)^{F(F-1)/2}=\ii^F,
\]
where the last equality holds because $F$ is even in the relevant cases.
With the use of the relations \eqref{eq:ch4-29-general-lorentz} to \eqref{eq:ch4-36-general-wlc} the proof goes just as
before.  We summarize:

\begin{theorem}[4-7 \emph{(PCT Theorem for General Spin)}]
\label{thm:ch4-7-general-spin-pct}

Let $\varphi_\mu\ldots\psi_\nu$ be spinor fields satisfying axioms I and II but
not necessarily III (LC).  If the $\PCT$ condition
\begin{equation*}
  \bra{\Omega}\varphi_\mu(x_1)\cdots\psi_\nu(x_n)\ket{\Omega}
  =\ii^F(-1)^J
  \bra{\Omega}\psi_\nu(-x_n)\cdots\varphi_\mu(-x_1)\ket{\Omega}
\end{equation*}
holds for all $x_1\ldots x_n$, then for every $x_1\ldots x_n$ such that
$x_1-x_2,\ldots,x_{n-1}-x_n$ is a Jost point, the WLC condition
\begin{equation*}
  \bra{\Omega}\varphi_\mu(x_1)\cdots\psi_\nu(x_n)\ket{\Omega}
  =\ii^F
  \bra{\Omega}\psi_\nu(x_n)\cdots\varphi_\mu(x_1)\ket{\Omega}
\end{equation*}
holds.  Conversely, if the WLC condition holds in a (real) neighborhood of a
Jost point, then the PCT conditions holds everywhere.

The normal commutation relations among the $\varphi\ldots\psi$ imply the WLC
conditions, so every field theory of fields with normal commutation relations
has $\PCT$ symmetry.
\end{theorem}

We shall see in the next section that if fields have “abnormal” commutation
relations there is also a symmetry $\Theta$ in the theory but it differs from
the definition \eqref{eq:ch1-general-spinor-pct-state} by an extra factor $\pm1$.  This is discussed at the end
of the next section.

It should be emphasized that while the assumption of local commutativity
implies that a field theory has $\PCT$ symmetry, only weak local commutativity
is necessary for this conclusion.  This is important in principle because it
is relatively straightforward to construct a field theory which satisfies weak
local commutativity (but not local commutativity) and has a non-trivial
$S$-matrix.  Such a theory will have $\PCT$ as a symmetry.  The $\PCT$ symmetry
observed in nature therefore is not a strong support for the hypothesis of
local commutativity.

\subsection{Spin and Statistics}
\label{sec:ch4-4-spin-statistics}

All experimental evidence indicates that systems with integer spin obey the
laws of Bose--Einstein statistics, and systems with half-odd integer spin those
of Fermi--Dirac statistics.  Although there are perfectly respectable laws of
statistics different from either Bose--Einstein or Fermi--Dirac, so far no
system has been seen to follow them \cite{ch4-28}.

A natural way to arrive at Bose--Einstein statistics is to describe the system
in question by a field which commutes for space-like separations, while the
analogous way for Fermi--Dirac statistics is to use a field which anti-commutes
for space-like separations.  The theorem on the connection of spin with
statistics or, as we shall say for brevity, the spin-statistics theorem, is an
assertion that in quantum field theory a non-trivial integer spin field cannot
have an anti-commutator vanishing for space-like separations, and a non-trivial
half-odd integer spin field cannot have a commutator vanishing for space-like
separations.  If one puts aside the possibility of laws of statistics other
than Bose--Einstein or Fermi--Dirac, the spin statistics theorem then accounts
for the experimental results.

When one turns from the commutation relations for a given field to those
between different fields the situation becomes more complicated.  It turns
out that “abnormal” commutation relations in which two integer spin fields, or
an integer spin and a half-odd integer spin field anti-commute, or two half-odd
integer spin fields commute, can be realized but, in general, the resulting
theories possess special symmetries.  By virtue of these symmetries it turns
out that there always exists another set of fields, satisfying normal
commutation relations and related to the original fields by a so-called Klein
transformation.  The original theory can equally well be regarded as a theory
of this set of fields.  In this sense, a theory with abnormal commutation
relations is a special case of a theory with normal commutation relations, one
which possesses a set of symmetries.

We shall prove these statements in turn.  Following Dell’Antonio we first
dispose of the possibility that a field component $\varphi$ has different
commutation relations with a field component $\psi$ and its adjoint $\psi^\dagger$.

\begin{theorem}[4-8]\label{thm:ch4-8}

If in a field theory we have
\begin{equation}
  [\varphi(x),\psi(y)]_{\pm}=0
  \qquad\text{for }(x-y)^2>0,
  \qquad\text{and the opposite}
  \qquad
  [\varphi(x),\psi^\dagger(y)]_{\mp}=0
  \tag{4-37}\label{eq:ch4-37-commutation-opposite}
\end{equation}
then either $\varphi$ or $\psi$ vanishes.
\end{theorem}

\begin{proof}

If $f$ and $g$ are two test functions of compact support we have
\begin{equation}
  \bra{\Omega}\varphi^\dagger(f)\psi^\dagger(g)\psi(g)\varphi(f)\ket{\Omega}
  =\bigl\lVert\psi(g)\varphi(f)\ket{\Omega}\bigr\rVert^2\geq0.
  \tag{4-38}\label{eq:ch4-38-positive-norm}
\end{equation}

If the supports of $f$ and $g$ are space-like separated, the assumed
commutation relations \eqref{eq:ch4-37-commutation-opposite} imply that the left-hand side of \eqref{eq:ch4-38-positive-norm} is
\begin{equation}
  -\bra{\Omega}\psi^\dagger(g)\psi(g)\varphi(f)\varphi^\dagger(f)\ket{\Omega}.
  \tag{4-39}\label{eq:ch4-39-reordered-negative}
\end{equation}

By the cluster decomposition property, if the support of $g$ runs to infinity
in a space-like direction, \eqref{eq:ch4-39-reordered-negative} approaches
\begin{equation*}
  -\bra{\Omega}\psi^\dagger(g)\psi(g)\ket{\Omega}
   \bra{\Omega}\varphi(f)\varphi^\dagger(f)\ket{\Omega}
  =-\bigl\lVert\psi(g)\ket{\Omega}\bigr\rVert^2
   \bigl\lVert\varphi^\dagger(f)\ket{\Omega}\bigr\rVert^2
\end{equation*}
and this is not positive.  Therefore, comparing with \eqref{eq:ch4-38-positive-norm} we see that zero
is the only consistent value of the limit, and either
$\psi(g)\ket{\Omega}=0$ or $\varphi^\dagger(f)\ket{\Omega}=0$.  By Theorem \ref{thm:ch4-3-separating}
this implies either $\psi(g)=0$ or $\varphi^\dagger(f)=0$.  If $\psi\ne0$
there exists a test function of compact support $g$ for which $\psi(g)\ne0$.
Then, for all test functions $f$ of compact support, $\varphi^\dagger(f)=0$.
Since the test functions of compact support are dense in $\mathcal F$, this
implies $\varphi=0$.  Similarly, if $\varphi\ne0$, then $\psi=0$.
\end{proof}

\begin{corollary*}

In a field theory there can be no non-zero field satisfying
\begin{equation}
  [\varphi(x),\varphi(y)]_{\pm}=0
  \tag{4-40}\label{eq:ch4-40-no-nonzero-field-first}
\end{equation}
and
\begin{equation*}
  [\varphi(x),\varphi^\dagger(y)]_{\mp}=0,
  \qquad\text{for all }(x-y)^2>0.
\end{equation*}
\end{corollary*}

Next we prove the proper spin-statistics theorem.  For clarity, we first deal
with a scalar field and then describe the modifications necessary for arbitrary
spin.

\begin{theorem}[4-9 \emph{(Spin-Statistics Theorem for a Scalar Field)}]
\label{thm:ch4-9-scalar-spin-statistics}

Let $\varphi$ be a scalar field.  Suppose that
\begin{equation}
  [\varphi(x),\varphi^\dagger(y)]_{\pm}=0
  \qquad\text{for }(x-y)^2>0.
  \tag{4-41}\label{eq:ch4-41-wrong-scalar-statistics}
\end{equation}

Then $\varphi(x)\ket{\Omega}=0=\varphi^\dagger(x)\ket{\Omega}$.  In a field
theory in which $\varphi$ and $\varphi^\dagger$ commute or anti-commute with
all other fields this implies $\varphi=\varphi^\dagger=0$.
\end{theorem}

\begin{proof}

The hypothesis \eqref{eq:ch4-41-wrong-scalar-statistics} of the “wrong” connection of spin with statistics
implies, if $(x-y)^2>0$,
\begin{equation}
  \bra{\Omega}\varphi(x)\varphi^\dagger(y)\ket{\Omega}
  +\bra{\Omega}\varphi^\dagger(y)\varphi(x)\ket{\Omega}=0.
  \tag{4-42}\label{eq:ch4-42-scalar-vacuum-relation}
\end{equation}

Each of the vacuum expectation values in this equation depends only on
$(x-y)$, and is the boundary value of a holomorphic function
\begin{equation}
\begin{aligned}
  \bra{\Omega}\varphi(x)\varphi^\dagger(y)\ket{\Omega}
  &=\lim_{\substack{\eta\to0\\\eta\in V_+}}
    W(x-y-\ii\eta),\\
  \bra{\Omega}\varphi^\dagger(y)\varphi(x)\ket{\Omega}
  &=\lim_{\substack{\eta\to0\\\eta\in V_+}}
  \widehat W(y-x-\ii\eta).
\end{aligned}
  \tag{4-43}\label{eq:ch4-43-holomorphic-boundaries}
\end{equation}
where $W$ and $\widehat W$ are holomorphic for $\eta\in V_+$.  Because $W$
and $\widehat W$ are invariant under the restricted Lorentz group, Theorem
\ref{thm:ch2-11} implies they are invariant under the proper complex Lorentz group, and
are holomorphic in the extended tube $\mathcal T'_1$ in the variable
$\zeta=(x-y)-\ii\eta$, which includes all points where $\zeta$ is real and
space-like.  Thus, \eqref{eq:ch4-42-scalar-vacuum-relation} is a relation between holomorphic functions
\begin{equation}
  W(\zeta)+\widehat W(-\zeta)=0,
  \tag{4-44}\label{eq:ch4-44-holomorphic-sum}
\end{equation}
which holds at an open set of real points of the extended tube and therefore
throughout the extended tube.  Using the invariance of $\widehat W$ under the
complex proper Lorentz transformation $\Lambda=-1$,
\begin{equation}
  \widehat W(\zeta)=\widehat W(-\zeta),
  \tag{4-45}\label{eq:ch4-45-hat-w-invariance}
\end{equation}
and we get from \eqref{eq:ch4-44-holomorphic-sum}
\begin{equation}
  W(\zeta)+\widehat W(\zeta)=0,
  \tag{4-46}\label{eq:ch4-46-holomorphic-zero}
\end{equation}
valid for $\zeta\in\mathcal T'_1$.

Passing to the boundary values in \eqref{eq:ch4-46-holomorphic-zero} as $\eta\to0$, $\eta\in V_+$, we
obtain for all $x$ and $y$ the distribution relation
\begin{equation}
  \bra{\Omega}\varphi(x)\varphi^\dagger(y)\ket{\Omega}
  +\bra{\Omega}\varphi^\dagger(-y)\varphi(-x)\ket{\Omega}=0.
  \tag{4-47}\label{eq:ch4-47-distribution-relation}
\end{equation}

We claim that this equation \eqref{eq:ch4-47-distribution-relation} implies $\varphi(x)\ket{\Omega}=0$.  To
prove this, set $\widehat f(x)=f(-x)$.  Then, recalling that
$\varphi^\dagger(f)=[\varphi(f)]^\dagger$,
\begin{equation*}
\begin{aligned}
  \varphi(f)&=\int\dd^4x\,f(x)\varphi(x),\\
  \varphi(\widehat f)&=\int\dd^4x\,f(-x)\varphi(x)\\
    &=\int\dd^4x\,f(x)\varphi(-x),
\end{aligned}
\end{equation*}
we get
\begin{equation*}
\begin{aligned}
  \bigl\lVert\varphi^\dagger(f)\ket{\Omega}\bigr\rVert^2
  +\bigl\lVert\varphi(\widehat f)\ket{\Omega}\bigr\rVert^2
  &=\bra{\Omega}\varphi(f)\varphi^\dagger(f)\ket{\Omega}
    +\bra{\Omega}\varphi^\dagger(\widehat f)\varphi(\widehat f)\ket{\Omega}\\
  &=\int\dd^4x\,\dd^4y\,f(x)\overline{f(y)}
    \bigl[\bra{\Omega}\varphi(x)\varphi^\dagger(y)\ket{\Omega}
    +\bra{\Omega}\varphi^\dagger(-y)\varphi(-x)\ket{\Omega}\bigr]\\
  &=0.
\end{aligned}
\end{equation*}

Thus, for all test functions, $f$,
$\lVert\varphi(f)\ket{\Omega}\rVert
=\lVert\varphi^\dagger(f)\ket{\Omega}\rVert=0$, which yields
$\varphi(f)\ket{\Omega}=0$ and completes the proof of the first half of the
theorem.

In any field theory in which all the fields commute or anti-commute, Theorem
\ref{thm:ch4-3-separating} applies, and so $\varphi(f)\ket{\Omega}=0$ implies
$\varphi=\varphi^\dagger=0$.
\end{proof}

For fields of general spin an analogue of Theorem \ref{thm:ch4-9-scalar-spin-statistics} is

\begin{theorem}[4-10 \emph{(Spin-Statistics Theorem for General Spin)}]
\label{thm:ch4-10-general-spin-statistics}

For a general irreducible spinor field the “wrong” connection of spin with
statistics
\begin{equation}
\begin{aligned}
  [\varphi_\alpha(x),\varphi_\alpha^\dagger(y)]_+&=0
    &&\varphi\text{ of integer spin},\\
  [\varphi_\alpha(x),\varphi_\alpha^\dagger(y)]_-&=0
    &&\varphi\text{ of half-odd integer spin},\\[-2pt]
  &\hspace{2cm}\text{for }(x-y)^2>0
\end{aligned}
  \tag{4-48}\label{eq:ch4-48-general-spin-wrong-statistics}
\end{equation}
implies $\varphi_\alpha(x)\ket{\Omega}=0$.  In a field theory in which all
fields either commute or anti-commute this implies
$\varphi_\alpha=\varphi_\alpha^\dagger=0$.
\end{theorem}

\begin{proof}

Without confusion we may leave out the suffix $\alpha$.  The proof follows the
same lines as for a scalar field.  As the generalization of \eqref{eq:ch4-42-scalar-vacuum-relation}, the
relations \eqref{eq:ch4-48-general-spin-wrong-statistics} give directly
\begin{equation}
  \bra{\Omega}\varphi(x)\varphi^\dagger(y)\ket{\Omega}
  \mathbin{\pm}
  \bra{\Omega}\varphi^\dagger(y)\varphi(x)\ket{\Omega}=0,
  \qquad\text{for }(x-y)^2>0,
  \tag{4-49}\label{eq:ch4-49-general-spin-vacuum}
\end{equation}
the sign being the same as in \eqref{eq:ch4-48-general-spin-wrong-statistics}.  As in the proof of Theorem \ref{thm:ch4-9-scalar-spin-statistics}, there
are holomorphic functions $W$ and $\widehat W$ such that the vacuum expectation
values are given by \eqref{eq:ch4-43-holomorphic-boundaries}.  The generalization of \eqref{eq:ch4-44-holomorphic-sum} is
\begin{equation}
  W(\zeta)\mathbin{\pm}\widehat W(-\zeta)=0,
  \tag{4-50}\label{eq:ch4-50-general-spin-holomorphic}
\end{equation}
while that of \eqref{eq:ch4-45-hat-w-invariance} is
\begin{equation}
  \widehat W(\zeta)=(-1)^J\widehat W(-\zeta),
  \tag{4-51}\label{eq:ch4-51-general-spin-invariance}
\end{equation}
where $J$ is the total number of undotted indices in $\varphi\varphi^\dagger$.
Equation \eqref{eq:ch4-51-general-spin-invariance} is a consequence of the transformation law of $\widehat W$
under the group $SL(2,\C)\otimes SL(2,\C)$ and is a special case of \eqref{eq:ch4-32-general-sign-reversal}.
We remark that the number of undotted indices in $\varphi\varphi^\dagger$ is
the same as the total number of indices in $\varphi$, and so is even (odd) if
the field $\varphi$ is of integer (half-odd integer) spin.  Therefore, at
points of holomorphy,
\begin{equation}
  \widehat W(\zeta)=\mathbin{\pm}\widehat W(-\zeta),
  \tag{4-52}\label{eq:ch4-52-general-spin-holomorphic-invariance}
\end{equation}
the sign being the same as in \eqref{eq:ch4-48-general-spin-wrong-statistics}.  Combining \eqref{eq:ch4-52-general-spin-holomorphic-invariance} and \eqref{eq:ch4-50-general-spin-holomorphic}, we arrive
at \eqref{eq:ch4-46-holomorphic-zero}, in exactly the same form as before.  From here on, the argument is
the same as for Theorem \ref{thm:ch4-9-scalar-spin-statistics}.
\end{proof}

Now we turn to a discussion of the commutation relations between
different fields.  We begin with an example, which illustrates a number of
the principles involved in a simple way.

\par\medskip\noindent\textbf{Example 1: \emph{Anti-commuting Scalar and Spin $\frac12$ Field}}
\par\smallskip

Suppose $\varphi$ is a scalar field and $\psi$ is a spin $\frac12$ field, and
that $\varphi$ and $\psi$ anti-commute at space-like separations.  We can
define new fields $\varphi'$ and $\psi'$ by
\begin{equation*}
\begin{aligned}
  \varphi'(x)\ket{\Psi_1}&=\varphi(x)\ket{\Psi_1},
  &\psi'(x)\ket{\Psi_1}&=\psi(x)\ket{\Psi_1},
\end{aligned}
\end{equation*}
if $\ket{\Psi_1}$ is a vector in the domain of definition $D$ of $\varphi,\psi$
in the coherent subspace $\mathcal H_1$ of $\mathcal H$ corresponding to an
integer spin representation of $\mathcal P_+^\uparrow$, and
\begin{equation*}
\begin{aligned}
  \varphi'(x)\ket{\Psi_2}&=-\varphi(x)\ket{\Psi_2},
  &\psi'(x)\ket{\Psi_2}&=\psi(x)\ket{\Psi_2},
\end{aligned}
\end{equation*}
if $\ket{\Psi_2}$ is a vector in $D$ in the coherent subspace $\mathcal H_2$
of $\mathcal H$ corresponding to a half-odd integer spin representation of
$\mathcal P_+^\uparrow$.  The action of $\varphi',\psi'$ on linear
superpositions of states from $\mathcal H_1$ and $\mathcal H_2$ is given by
linearity.

A super-selection rule (the univalence super-selection rule) operates between
$\mathcal H_1$ and $\mathcal H_2$.  Recall from Section \ref{sec:ch1-super-selection-rules} that linear
superpositions of states from different coherent subspaces between which
there is a super-selection rule are not physically realizable.  In this case,
a state $\alpha\ket{\Psi_1}+\beta\ket{\Psi_2}$ becomes
$\alpha\ket{\Psi_1}-\beta\ket{\Psi_2}$ when it is rotated through an angle of
$2\pi$ around any axis.  This should belong to the same ray, and so that state
is not physically realizable unless $\alpha=0$ or $\beta=0$.

It is easy to see that the new fields, $\varphi',\psi'$ commute at space-like
separated points, that is, in contrast to $\varphi,\psi$ the new fields satisfy
the normal commutation relations.  The transformation from $\varphi,\psi$ to
$\varphi',\psi'$ is called a Klein transformation.  Unlike the case of the
symmetries we have discussed in Section \ref{sec:ch3-symmetries}, there is no unitary or anti-unitary
transformation $V$ such that
\begin{equation*}
  \varphi'=V\varphi V^{-1},
  \qquad
  \psi'=V\psi V^{-1},
\end{equation*}
because the existence of such a $V$ would imply
\begin{equation*}
  [\varphi'(x),\psi'(y)]_-=V[\varphi(x),\psi(y)]_-V^{-1}.
\end{equation*}

We can take two points of view with respect to the transformation
$\varphi,\psi\to\varphi',\psi'$.  In the first, it is simply “a change of
variables.”  The observables of the theory will be certain functions of the
$\varphi,\psi$, and $\varphi,\psi$ are functions of the $\varphi',\psi'$.
Hence the observables are certain (different) functions of the fields
$\varphi',\psi'$.  The preceding discussion merely says that by making a
change of variables one can express the observables in terms of a set of
fields having normal commutation relations.

In the second point of view, one makes a different correspondence between
operators and laboratory operations.  If a certain function,
say $F(\varphi,\psi)$, is observable and corresponds to certain well-defined
measurements in the laboratory, then the Klein transformation is regarded as
leading to a new theory in which $F(\varphi',\psi')$ corresponds to that very
set of measurements.  It is natural to ask whether the two theories predict
the same result for all experiments.  Whether this is so only a detailed
analysis of the measurement theory of the system can decide.  The more
operators that are required to be observable, the less chance there is that the
two theories are physically equivalent.  In our example we have
\begin{equation*}
  \bra{\Psi_2}\varphi(x)\ket{\Psi_2}
  =-\bra{\Psi_2}\varphi'(x)\ket{\Psi_2},
\end{equation*}
so if $\varphi(x)$ is observable then the two theories are different.  At the
other extreme, where only the results of scattering experiments are
observable, the theories predict the same results; this is discussed below in
connection with the $\PCT$ theorem.

The example just considered has one feature not characteristic of the general
situation for abnormal commutation relations; the existence of the orthogonal
subspaces $\mathcal H_1$ and $\mathcal H_2$ invariant under $U(a,A)$ is already
a consequence of the univalence super-selection rule.  In general, the
existence, orthogonality, and invariance of the analogous subspaces is a
consequence of the abnormal commutation relations themselves.  Since several
new points of principle arise in this situation, we illustrate it by means of
a second simple example.

\par\medskip\noindent\textbf{Example 2: \emph{A Scalar and Two Spin $\frac12$ Fields}}
\par\smallskip

Consider a theory of one hermitian scalar field $\varphi$ and two hermitian
half-odd integer spin fields $\psi_1$ and $\psi_2$.  For simplicity the spinor
indices on $\psi_1$ and $\psi_2$ will be suppressed.  Suppose the normal
commutation relations
\begin{equation}
  [\psi_1(x),\psi_2(y)]_+=0=[\varphi(x),\psi_2(y)]_-
  \qquad (x-y)^2>0
  \tag{4-53}\label{eq:ch4-53-example-2-normal}
\end{equation}
hold, together with the “abnormal” relation
\begin{equation}
  [\varphi(x),\psi_1(y)]_+=0,
  \qquad (x-y)^2>0.
  \tag{4-54}\label{eq:ch4-54-example-2-abnormal}
\end{equation}

These relations imply the vanishing of certain vacuum expectation values, as
will now be shown.  For later application the result will be stated and proved
for general field theories.

\begin{theorem}[4-11]\label{thm:ch4-11}

In any local field theory, if $M(x_1,\ldots,x_j)$ and $N(y_1,\ldots,y_k)$ are
two monomials in the field components which anti-commute if $(x)$ and
$(y)$ are space-like separated sets of points, then either
$\bra{\Omega}M\ket{\Omega}=0$, or $\bra{\Omega}N\ket{\Omega}=0$.
\end{theorem}

\begin{proof}

Consider
\begin{equation}
\begin{aligned}
  &\bra{\Omega}M(x_1,\ldots,x_j)N(y_1+a,\ldots,y_k+a)\ket{\Omega}\\
  &\quad=-\bra{\Omega}N(y_1+a,\ldots,y_k+a)M(x_1,\ldots,x_j)\ket{\Omega}
\end{aligned}
  \tag{4-55}\label{eq:ch4-55-theorem-4-11-cluster}
\end{equation}
if $a$ is a large space-like vector.  The limit of this equation as
$a\to\infty$ is, by the cluster decomposition theorem (Theorem \ref{thm:ch3-4-cluster-decomposition}),
\begin{equation*}
  \bra{\Omega}M\ket{\Omega}\bra{\Omega}N\ket{\Omega}
  =-\bra{\Omega}N\ket{\Omega}\bra{\Omega}M\ket{\Omega}.
\end{equation*}

Therefore, one or the other of $\bra{\Omega}M\ket{\Omega}$ and
$\bra{\Omega}N\ket{\Omega}$ must vanish.
\end{proof}

Returning again to the model, let us suppose that for some odd $n$,
$\bra{\Omega}\varphi(x_1)\cdots\varphi(x_n)\ket{\Omega}\ne0$.  The
alternative possibility can, of course, also be analyzed, but this case
suffices to illustrate our point.  Putting
\begin{equation*}
\begin{aligned}
  M&=\varphi(x_1)\cdots\varphi(x_n), & n&\ \text{odd}.
\end{aligned}
\end{equation*}

\begin{equation*}
\begin{aligned}
  N={}&\varphi(y_1)\cdots\varphi(y_j)
  \psi_1(y_{j+1})\cdots\psi_1(y_{j+k})\\
  &\phantom{={}}\times\psi_2(y_{j+k+1})\cdots\psi_2(y_{j+k+\ell}),
  & k&\ \text{odd},\quad \ell\ \text{odd}.
\end{aligned}
\end{equation*}
and applying Theorem~\ref{thm:ch4-11}, since $M$ and $N$ anti-commute we get for all $j$
\begin{equation}
\begin{aligned}
  \bra{\Omega}\varphi(y_1)\cdots\varphi(y_j)
  \psi_1(y_{j+1})\cdots\psi_1(y_{j+k})
  \psi_2(y_{j+k+1})\cdots\psi_2(y_{j+k+\ell})\ket{\Omega}
  &=0,\\[-2pt]
  k&\ \text{odd},\quad \ell\ \text{odd}.
\end{aligned}
  \tag{4-56}\label{eq:ch4-56-example-2-vacuum-zero}
\end{equation}

The expression already vanishes if $k+\ell$ is odd [see the remark after
\eqref{eq:ch4-30-general-boundary}].  Clearly this and \eqref{eq:ch4-56-example-2-vacuum-zero} together imply that the vacuum expectation
values, and therefore the theory, are invariant under the group of four
transformations given by $(\psi_1,\psi_2)\to(\pm\psi_1,\pm\psi_2)$.  This
symmetry gives rise to a conservation law, which will be called an
\emph{even-odd rule}.

The Hilbert space $\mathcal H$ splits into a sum of four orthogonal subspaces
$\mathcal H_{+,+},\mathcal H_{+,-},\mathcal H_{-,+},\mathcal H_{-,-}$, spanned
respectively by states
\begin{equation*}
  \varphi(x_1)\cdots\varphi(x_j)\psi_1(y_{j+1})\cdots\psi_1(y_{j+k})
  \psi_2(y_{j+k+1})\cdots\psi_2(y_{j+k+\ell})\ket{\Omega}
\end{equation*}
with $(k,\ell)=(\text{even},\text{even}), (\text{even},\text{odd}),
(\text{odd},\text{even}),$ and $(\text{odd},\text{odd})$.  The orthogonality
of the four subspaces is guaranteed by \eqref{eq:ch4-56-example-2-vacuum-zero} together with the univalence
super-selection rule.  Clearly, they are invariant under the transformations
of the Poincaré group.

We now define the Klein transformation for this case.  It leaves $\psi_1$ and
$\psi_2$ invariant and replaces $\varphi$ by a new field $\varphi'$, defined
as $\varphi$ on those vectors of the domain of $\varphi$ which are in
$\mathcal H_{+,\pm}$, $-\varphi$ on those which are in
$\mathcal H_{-,\pm}$, and by linearity elsewhere.  It can be seen that
$\varphi'$ is hermitian.  Then $\varphi'$, together with $\psi'_1=\psi_1$ and
$\psi'_2=\psi_2$, satisfy the normal commutation relations.

Analogues of the remarks made on the first example apply here.  In addition,
there are the following new features.  The even-odd rule defines a
conservation law for the system; does it define a super-selection rule?  Again,
an analysis of the observables is necessary to answer such questions.  The
new theory is equivalent to the old in the first sense discussed above: there
exists an irreducible set of field operators $(\varphi',\psi'_1,\psi'_2)$ which
have normal commutation relations.

It may be, of course, that interesting quantities of the theory are more
“simply” expressible in terms of one set of fields than the other.  In this
connection there is a good reason for preferring the transformed fields
$(\varphi',\psi'_1,\psi'_2)$, for the hermitian quantities $\varphi(x)$ and
$\psi_1(y)\psi_2(z)+\psi_2(z)\psi_1(y)$ do not commute at large space-like
distances.  This seems unnatural in a field theory; it means that any function
of the fields $(\varphi,\psi_1,\psi_2)$ which corresponds to a local
measurement of the system must be a rather complicated function of the old
fields.  It might be a simple function of the new fields, which are not beset
by lack of commutativity at space-like separations.  Our discussion of the
second example is now complete.

We now deal with the case of a general set of local fields, following the
treatment given by Araki.  The essential steps are to show that abnormal
commutation relations give rise to even-odd rules and then to show that this
makes possible the definition of the required Klein transformations.  We shall
suppose that there are $n$ fields $\varphi_1,\ldots,\varphi_n$ and that all
components of $\varphi_j$ satisfy the same commutation (or anti-commutation)
relations with all components of $\varphi_k$.  We shall assume the adjoints
$\varphi_{j\alpha}^\dagger$ of the components $\varphi_{j\alpha}$ of
$\varphi_j$ occur somewhere in the list of $\varphi_{k\beta}$.  Our ultimate
objective is given in the following theorem.

\begin{theorem}[4-12]\label{thm:ch4-12}

In any field theory with abnormal commutation relations there always exists an
irreducible set of fields with normal commutation relations obtained from the
original set by a Klein transformation.
\end{theorem}

\par\medskip\noindent The proof is rather involved, and will be carried out in
a number of stages.

The first step is to generalize the notion of an even-odd rule.  A theory is
said to have an \emph{even-odd rule} for a set $\alpha$ of fields if every
vacuum expectation value with an odd number of fields from $\alpha$ vanishes.
In such a case the two spaces $\mathcal H_o$ and $\mathcal H_e$ generated from
the vacuum
by polynomials which are odd or even, respectively, in the fields belonging
to $\alpha$ (and arbitrary in the rest of the fields), are orthogonal to each
other, and are obviously invariant under $\mathcal P_+^\uparrow$.  Hence the
operator $p(\alpha)$, defined to be $1$ on $\mathcal H_e$ and $-1$ on
$\mathcal H_o$, and by linearity elsewhere, commutes with the representation
$U(a,\Lambda)$ of $\mathcal P_+^\uparrow$.

Suppose we define a Klein transformation to new fields by multiplying all the
fields in another subset $\beta$ of the fields by $p(\alpha)$:
\begin{equation}
\begin{aligned}
  \varphi'_j&=p(\alpha)\varphi_j, &&\varphi_j\in\beta,\\
  \varphi'_j&=\varphi_j, &&\varphi_j\notin\beta.
\end{aligned}
  \tag{4-57}\label{eq:ch4-57-klein-general}
\end{equation}

Then $\varphi'$ again are fields, in general with different commutation
relations.  The set $\alpha$ in \eqref{eq:ch4-57-klein-general} (or several such sets) is determined by
the abnormal commutation relations, i.e., by applying Theorem \ref{thm:ch4-11}.  The set
$\beta$ is free to be chosen appropriately so that the commutation relations
are changed to the normal ones.  In the first example above, $\alpha$ is the
set consisting of the spin $\frac12$ field $\psi$, and the even-odd rule is
that given by the univalence super-selection rule.  In the second example
there are two possible $\alpha$'s, the set consisting of $\psi_1$ or the set
consisting of $\psi_2$.  Our choice of the set $\beta$ was, in both examples,
the set consisting of the single field $\varphi$.

The commutation relations of $\varphi'_j,\varphi'_k$ are different from those
of $\varphi_j,\varphi_k$ if one field is in $\alpha$ and the other is in
$\beta$, and both fields are not in both sets.  Otherwise the relations are
not altered.  To see this we must consider in turn all the possible cases, of
which there are sixteen, enumerated according as $\varphi_1$ or $\varphi_2$
is or is not in the set $\alpha$ or $\beta$.  (For simplicity, $j$ and $k$
are replaced by 1 and 2 temporarily.)  If neither $\varphi_1$ nor
$\varphi_2$ lies in the set $\beta$, then $\varphi'_1=\varphi_1$ and
$\varphi'_2=\varphi_2$, and so the relations are unchanged.  This takes care
of four cases, labeled by
$(\varphi_1\in\alpha,\ \varphi_1\notin\beta;\ \varphi_2\in\alpha,\ \varphi_2\notin\beta)$;
$(\varphi_1\notin\alpha,\ \varphi_1\notin\beta;\ \varphi_2\in\alpha,\ \varphi_2\notin\beta)$;
$(\varphi_1\in\alpha,\ \varphi_1\notin\beta;\ \varphi_2\notin\alpha,\ \varphi_2\notin\beta)$;
$(\varphi_1\notin\alpha,\ \varphi_1\notin\beta;\ \varphi_2\notin\alpha,\ \varphi_2\notin\beta)$.

If neither $\varphi_1$ nor $\varphi_2$ is in $\alpha$ then both fields commute
with $p(\alpha)$, and the commutation rule is unchanged.  This accounts for
three more cases:
$(\varphi_1\notin\alpha,\ \varphi_1\notin\beta;\ \varphi_2\notin\alpha,\ \varphi_2\in\beta)$;
$(\varphi_1\notin\alpha,\ \varphi_1\in\beta;\ \varphi_2\notin\alpha,\ \varphi_2\in\beta)$;
and $(\varphi_1\notin\alpha,\ \varphi_1\in\beta;\ \varphi_2\notin\alpha,\ \varphi_2\notin\beta)$.

It is also easy to see that if one of the fields is in neither set, the Klein
transformation does not alter the commutator.  This accounts for
$(\varphi_1\notin\alpha,\ \varphi_1\notin\beta;\ \varphi_2\in\alpha,\ \varphi_2\in\beta)$;
$(\varphi_1\in\alpha,\ \varphi_1\in\beta;\ \varphi_2\notin\alpha,\ \varphi_2\notin\beta)$.

If both $\varphi_1$ and $\varphi_2$ lie in both sets, then
\[
  \varphi'_1\varphi'_2=p(\alpha)\varphi_1p(\alpha)\varphi_2
  =-p(\alpha)^2\varphi_1\varphi_2=-\varphi_1\varphi_2
\]
and $\varphi'_2\varphi'_1=-\varphi_2\varphi_1$.  Therefore the vanishing of
$[\varphi_1,\varphi_2]_\pm$ is unchanged by \eqref{eq:ch4-57-klein-general}.  This disposes of
$(\varphi_1\in\alpha,\ \varphi_1\in\beta;\ \varphi_2\in\alpha,\ \varphi_2\in\beta)$.
The remaining six possibilities satisfy the condition that one of
$\varphi_1,\varphi_2$ lies in $\alpha$ and the other in $\beta$.  They are
$(\varphi_1\notin\alpha,\ \varphi_1\in\beta;\ \varphi_2\in\alpha,\ \varphi_2\notin\beta)$;
$(\varphi_1\notin\alpha,\ \varphi_1\in\beta;\ \varphi_2\in\alpha,\ \varphi_2\in\beta)$;
$(\varphi_1\in\alpha,\ \varphi_1\notin\beta;\ \varphi_2\notin\alpha,\ \varphi_2\in\beta)$;
$(\varphi_1\in\alpha,\ \varphi_1\notin\beta;\ \varphi_2\in\alpha,\ \varphi_2\in\beta)$;
$(\varphi_1\in\alpha,\ \varphi_1\in\beta;\ \varphi_2\notin\alpha,\ \varphi_2\in\beta)$;
$(\varphi_1\in\alpha,\ \varphi_1\in\beta;\ \varphi_2\in\alpha,\ \varphi_2\notin\beta)$.

It is left to the reader to check that in these six cases,
$[\varphi_1,\varphi_2]_\pm=0$ changes into
$[\varphi'_1,\varphi'_2]_\mp=0$, using \eqref{eq:ch4-57-klein-general} and the fact that if
$\varphi_j\in\alpha$, $p(\alpha)$ anti-commutes with $\varphi_j$, whereas if
$\varphi_j\notin\alpha$, $p(\alpha)$ and $\varphi_j$ commute.  We therefore
have the result: the transformation \eqref{eq:ch4-57-klein-general} alters the commutation relations
if one field is in $\alpha$ and the other is in $\beta$ and both fields are
not in both sets; otherwise the relations are unchanged.

We want to transform the commutation relations to the normal ones, and so the
matrix $\sigma_{ij}$ defined as follows is of interest:
\begin{equation}
\begin{aligned}
  \sigma_{ij}&=0 &&\text{if }\varphi_i,\varphi_j\text{ have normal commutation relations},\\
  \sigma_{ij}&=1 &&\text{otherwise.}
\end{aligned}
  \tag{4-58}\label{eq:ch4-58-sigma-definition}
\end{equation}

Obviously $\sigma_{ij}=\sigma_{ji}$; Theorem \ref{thm:ch4-10-general-spin-statistics} asserts that a field always
has normal commutation relations with itself, which means that $\sigma_{ii}=0$.

In order to determine whether we have an even-odd rule we examine the
commutation relations of monomials $M,N,\ldots$ of the fields, and we can
introduce the concept of \emph{normal commutation relations between
monomials}: two monomials $M$ and $N$ are said to have “normal” commutation
relations if they commute (anti-commute) just as if all the fields in them had
normal commutation relations.  Otherwise the monomials are said to have
“abnormal” commutation relations.  Two monomials may have normal commutation
relations although their constituent fields have not.  It is clear that the
commutation relations of one monomial with another depends only on the
parities of the powers of each field in them.\footnote{%
Two numbers are said to have the same \emph{parity} if they are both even,
or both odd.}

In other words, two monomials, $M_1$ and $M_2$, have the same relations with
any other monomial if $M_1$ and $M_2$ belong to the same equivalence class
defined as follows: two monomials are equivalent if the number of times each
field $\varphi_i$ occurs has the same parity in both.  Each equivalence class
contains a monomial of least degree, in which a given field $\varphi_j$ either
occurs once or not at all (the order in which the fields are multiplied does
not matter for our purposes).  Thus, there is a 1:1 correspondence between
the equivalence classes and subsets of the fields
$\varphi_1,\ldots,\varphi_n$.  We shall label each equivalence class by a
vector $s=(s_1,\ldots,s_n)$ defined as follows:
\begin{equation*}
  s_j=\begin{cases}0&\text{if }\varphi_j\text{ appears an even number of times},\\
                    1&\text{if }\varphi_j\text{ appears an odd number of times}
       \end{cases}
\end{equation*}
in the monomials in the equivalence class $s$.  If $M$ is a monomial we will
use the symbol $s(M)$ to denote the equivalence class containing
$M$ or the vector which labels it.  To any set $\beta$ of fields we can
associate the vector $t(\beta)=(t_1(\beta),\ldots,t_n(\beta))$, which is defined
by
\begin{equation*}
  t_j(\beta)=\begin{cases}
    0&\text{if }\varphi_j\notin\beta,\\
    1&\text{if }\varphi_j\in\beta.
  \end{cases}
\end{equation*}

If $\beta$ is the set $(\varphi_1,\ldots,\varphi_k)$, then clearly $t(\beta)$
is the equivalence class containing the monomial
$\varphi_1\cdots\varphi_k$.

We shall define addition of two vectors $s,t$ by
\begin{equation*}
  (s+t)_j=s_j+t_j\pmod 2,
\end{equation*}
that is, we use the addition rule for components $0+0=1+1=0$;
$0+1=1+0=1$.  This is convenient, since we are only interested in the parity
of the numbers involved.  With this addition law the equivalence classes
$s,\ldots$ form a vector space $V$.  It is easy to derive the following two
formulas:
\begin{equation}
  \sum_j s_j(M)t_j(\beta)=
  \left(\begin{smallmatrix}0\\1\end{smallmatrix}\right)\pmod 2.
  \tag{4-59}\label{eq:ch4-59-parity-sum}
\end{equation}
if $M$ contains an $\left(\begin{smallmatrix}\text{even}\\\text{odd}\end{smallmatrix}\right)$
number of fields from $\beta$, and from \eqref{eq:ch4-58-sigma-definition},
\begin{equation}
  \sum_{i,j}s_i(M)\sigma_{ij}s_j(N)=
  \left(\begin{array}{ll}
    0&\text{if $M,N$ obey the normal commutation relations},\\
    1&\text{otherwise,}
  \end{array}\right)\pmod 2.
  \tag{4-60}\label{eq:ch4-60-normal-monomial-pairing}
\end{equation}
the sum again being mod 2.

We now have most of the apparatus necessary to prove the theorem.  The effect
of the transformation \eqref{eq:ch4-57-klein-general} on $\sigma_{ij}$ can be calculated explicitly.
If $i$ and $j$ are such that $\varphi_i$ is in one of the sets $\alpha,\beta$,
and $\varphi_j$ is in the other, then from the definition of $t(\alpha)$,
$t(\beta)$
\begin{equation}
  \text{either }t_i(\alpha)=t_j(\beta)=1
  \qquad\text{or}\qquad t_i(\beta)=t_j(\alpha)=1.
  \tag{4-61}\label{eq:ch4-61-set-membership}
\end{equation}

We have remarked that $\sigma_{ij}$ is changed by the Klein transformation
\eqref{eq:ch4-57-klein-general} if and only if one and only one of \eqref{eq:ch4-61-set-membership} holds.  Thus the effect of
\eqref{eq:ch4-57-klein-general} on $\sigma$ is
\begin{equation}
  \sigma_{ij}\longrightarrow\sigma_{ij}
  +t_i(\alpha)t_j(\beta)+t_j(\alpha)t_i(\beta)\pmod 2.
  \tag{4-62}\label{eq:ch4-62-sigma-update}
\end{equation}

We show that $\sigma$ can be reduced to 0 by a series of transformations
\eqref{eq:ch4-57-klein-general} by proving that $\sigma$ has the form given by the following theorem.

\begin{theorem}[4-13]\label{thm:ch4-13}

If $\sigma_{ij}=\left(\begin{smallmatrix}0\\1\end{smallmatrix}\right)$ if
$\varphi_i$ and $\varphi_j$ have
$\left(\begin{smallmatrix}\text{normal}\\\text{abnormal}\end{smallmatrix}\right)$
commutation relations, then
\begin{equation}
  \sigma_{ij}=\sum_{k=1}^{N}
  \bigl[t_i(\alpha_k)s_j^{(k)}+s_i^{(k)}t_j(\alpha_k)\bigr]\pmod 2,
  \tag{4-63}\label{eq:ch4-63-sigma-decomposition}
\end{equation}
where the $\alpha_k$ are sets having an even-odd rule, and the $s^{(k)}$,
$k=1,\ldots,N$ are some vectors in $V$.
\end{theorem}

\begin{proof}[Proof of Theorem \ref{thm:ch4-12}]

Assuming the result of Theorem \ref{thm:ch4-13} we can proceed as follows.  Choose the set
$\beta_k$ to consist of the fields in the minimal monomial in $s^{(k)}$ of
minimal degree in each of the fields [i.e., $s^{(k)}=t(\beta_k)$].  If we
apply \eqref{eq:ch4-57-klein-general} $N$ times using the $\alpha_k,\beta_k$, then $\sigma_{ij}$ is
transformed to
\begin{equation*}
  [\sigma_{ij}+\sigma_{ij}]\pmod 2=0\pmod 2.
\end{equation*}
\end{proof}

\begin{proof}[Proof of Theorem \ref{thm:ch4-13}]

To prove Theorem \ref{thm:ch4-13} we must have a criterion for deciding whether a given
set $\alpha$ has an even-odd law.  For this we introduce the subset $\Gamma$
of $V$ defined by: $s\in\Gamma$ if there exists a monomial $M$ in the
equivalence class $s$ such that $\bra{\Omega}M\ket{\Omega}\ne0$.  We can
define $\overline\Gamma$ to be the set obtained by adjoining to $\Gamma$ all
sums of vectors in $\Gamma$.  [For example, if $\Gamma$ contains $(0,1)$ and
$(1,1)$, then $\overline\Gamma$ contains $(0,1)$, $(1,1)$, and
$(0,1)+(1,1)=(1,0)$.]  The set $\Gamma$ (or $\overline\Gamma$) is just what
is wanted in order to determine the even-odd rules of the theory, as will now
be explained.

Our assertion is the following: if
\begin{equation}
  \sum_j s_jt_j(\alpha)=0\pmod 2
  \qquad\text{for all }s\in\Gamma
  \tag{4-64}\label{eq:ch4-64-even-odd-criterion}
\end{equation}

(and incidentally, therefore, for all $s\in\overline\Gamma$), then there is an
even-odd rule for the set $\alpha$.  For the proof, suppose, if possible, that
$M$ is a monomial containing an odd number of fields from $\alpha$ (and any
number of other fields) and is such that $\bra{\Omega}M\ket{\Omega}\ne0$.
Let $s(M)$ be the equivalence class of $M$.  By the definition of $\Gamma$,
$s(M)\in\Gamma$, and by \eqref{eq:ch4-59-parity-sum}, $\sum_js_j(M)t_j(\alpha)=1\pmod2$,
contradicting the given property.

We next assert that if $s(M),s(N)$ are two vectors in $\Gamma$, then $M,N$
have normal commutation relations.  For $\bra{\Omega}M\ket{\Omega}\ne0$,
$\bra{\Omega}N\ket{\Omega}\ne0$ and therefore $M$ and $N$ commute to avoid
contradicting Theorem \ref{thm:ch4-11}.  Both $M$ and $N$ contain an even number of
half-odd integer
spin fields and so, if we had assumed that all fields $\varphi_j$ had normal
relations, $M$ and $N$ would commute.  Thus, $M$ and $N$ have normal
commutation relations.  Equivalent to this statement, according to \eqref{eq:ch4-60-normal-monomial-pairing}, is
the condition
\begin{equation}
  \sum_{i,j}s_i(M)\sigma_{ij}s_j(N)=0
  \qquad\text{if }s(M),s(N)\in\Gamma,
  \tag{4-65}\label{eq:ch4-65-gamma-orthogonality}
\end{equation}
which will play a significant role in the following.  This concludes the
preliminary discussion, and we now turn to the proof proper.

Let $e^{(1)},\ldots,e^{(m)}$ be a linearly independent set of vectors which
span $\overline\Gamma$, and $e^{(1)},\ldots,e^{(n)}$ a basis in $V$ which
includes $e^{(1)},\ldots,e^{(m)}$.

We can always find a (dual) linearly independent set $d^{(k)}$ such that
\begin{equation}
  \bigl(e^{(j)},d^{(k)}\bigr)=\sum_i e_i^{(j)}d_i^{(k)}=\delta_{jk}\pmod 2.
  \tag{4-66}\label{eq:ch4-66-dual-basis}
\end{equation}

We are here using the natural scalar product in $V$:
\begin{equation*}
  (s,t)=\sum_{i=1}^{n}s_it_i\pmod 2.
\end{equation*}

Note that this scalar product is not positive definite, so that mutual
orthogonality for a set of vectors does not imply linear independence.  For
example, $((1,1),(1,1))=1+1=0$ and so $(1,1)$ is orthogonal to itself.
However, this lack of positive definiteness does not prevent us from finding
the $d^{(k)}$.

We can regard $\sigma$ as an operator in $V$: $s\mapsto\sigma s$, given by
\begin{equation*}
  (\sigma s)_i=\sum_j\sigma_{ij}s_j\pmod 2.
\end{equation*}

Let us denote the matrix elements of $\sigma$ in the new coordinate system
$e^{(1)},\ldots,e^{(n)}$ by
\begin{equation}
  \sigma'_{ij}=\bigl(e^{(i)},\sigma e^{(j)}\bigr).
  \tag{4-67}\label{eq:ch4-67-new-coordinate-matrix}
\end{equation}

Then
\begin{equation}
  (t,\sigma s)=\sum_{i,k}(t,d^{(i)})\sigma'_{ik}(d^{(k)},s);
  \qquad s,t\in V.
  \tag{4-68}\label{eq:ch4-68-coordinate-form}
\end{equation}

We know that any monomial has normal relations with itself.  This applies to
the monomials defined by the vectors $e^{(i)}$, and so, using \eqref{eq:ch4-67-new-coordinate-matrix} we see
that $\sigma'_{ii}=0$, $i=1,2,\ldots,n$.  From the definition we easily see
that
\begin{equation*}
  \sigma'_{ij}=\sigma'_{ji}.
\end{equation*}

Because the vectors $e^{(1)},\ldots,e^{(m)}$ lie in $\overline\Gamma$,
$\sigma'_{ik}=0$ if $i,k\leq m$, as a consequence of \eqref{eq:ch4-65-gamma-orthogonality}.  Therefore
\begin{equation*}
  \sigma_{ij}=\sum_{k>m}
  \bigl(d_i^{(k)}s_j^{(k)}+s_i^{(k)}d_j^{(k)}\bigr),
\end{equation*}
where
\begin{equation*}
  s^{(k)}=\sum_{i<k}\sigma'_{ik}d^{(i)}.
\end{equation*}

We now assert that the vectors $d^{(k)}$ define sets $\alpha_k$ with an
even-odd rule, if $k>m$.  Let $\alpha_k$ consist of those fields present in
the monomial in the equivalence class defined by $d^{(k)}$ of minimal degree in
each of the fields, that is,
\begin{equation*}
\begin{aligned}
  d^{(k)}=t(\alpha_k),\qquad
  &\varphi_j\in\alpha_k &&\text{if }d_j^{(k)}=1,\\
  &\varphi_j\notin\alpha_k &&\text{if }d_j^{(k)}=0.
\end{aligned}
\end{equation*}

Then if $k>m$ we have $(d^{(k)},e^{(j)})=0$ for $j\leq m$, which gives
$\sum_i d_i^{(k)}s_i=0\pmod 2$ for all $s\in\Gamma$, since $e^{(j)}$,
$j=1,\ldots,m$, span $\Gamma$.  Therefore, by \eqref{eq:ch4-64-even-odd-criterion}, the set $\alpha_k$
defines an even-odd rule.
\end{proof}

We now return to the discussion of the $\PCT$ theorem, proved in the last
section to hold in any theory with the normal commutation relations.  What is
the situation when the commutation relations are abnormal?  The easiest way to
answer this question is to use the existence of the Klein transformation of
Theorem \ref{thm:ch4-12} from the given fields $\varphi_i$ to fields $\varphi'_i$
satisfying normal commutation relations.

Applying the $\PCT$ theorem to the new fields $\varphi'_i$ we see that there
exists a symmetry $\Theta$ in the theory which satisfies
\begin{equation}
  \Theta\varphi'_i(f)\Theta^{-1}
  =(-1)^J\left(\begin{smallmatrix}\ii\\1\end{smallmatrix}\right)
  \varphi_i^{\prime\dagger}(\widehat f),
  \qquad
  \Theta\ket{\Omega}=\ket{\Omega}.
  \tag{4-69}\label{eq:ch4-69-pct-abnormal}
\end{equation}
the upper case for half-odd integer spin fields, the lower case for integer
spin fields.  In \eqref{eq:ch4-69-pct-abnormal}, $\widehat f(x)=f(-x)$.  This will induce a
transformation of the original fields $\varphi_i$ which will, in general, not
be of the form \eqref{eq:ch4-69-pct-abnormal} with the primes removed.  To see how this comes about
consider a third example.

\par\medskip\noindent\textbf{Example 3: \emph{Anti-Commuting Hermitian Scalar Fields}}
\par\smallskip

Let $\varphi$ and $\psi$ be two anti-commuting hermitian scalar fields such
that
\begin{equation*}
  \bra{\Omega}\varphi(x)\psi(y)\ket{\Omega}\ne0.
\end{equation*}

We conclude from Theorem \ref{thm:ch4-11} that all vacuum expectation values containing an
odd number of fields vanish.  We can define a Klein transformation for this
case by $\varphi'=\varphi$, $\psi'=\psi$ on the subspace spanned by even
polynomials in the smeared fields applied to the vacuum and
